\documentclass{zhvt-classic}

% 使用本地源流明体（GenRyuMin2 TW）按文件路径加载，不依赖系统字体名
% 版框（内框／界栏／书口侧栏／鱼尾）改用与版心印章同一印泥棕色：
%   印章图片 yshs.png 实色 = RGB(135,68,52) = HTML 874434
\usepackage{xcolor}
\usepackage{tikz}
% sealbrown 双注册：xcolor 供 TikZ/pgf，L3 供 cls 版框
\definecolor{sealbrown}{HTML}{874434}
\ExplSyntaxOn
\color_set:nnn { sealbrown } { HTML } { 874434 }
\ExplSyntaxOff

\zhvtset{
  grid color       = sealbrown,
  main font        = GenRyuMin2TW-SB,
  main font path   = C:/Users/Administrator/texmf/fonts/opentype/public/GenRyuMin2TW/,
jiazhu font      = GenRyuMin2TW-M,
  jiazhu font path = C:/Users/Administrator/texmf/fonts/opentype/public/GenRyuMin2TW/,
  % 页码框容纳 4 字：正常页「九十七」+前置空 1 格=4 行恰好；
  % 100+ 页「一百零一」共 5 行超出，末字「一/二/三」被框裁掉（2026-10-04 实测），故放宽到 5
  shukou page number chars = 5,
  % 书口书名框高从默认 6 字减到 4 字，把省出的 2 字高拨给 chapter 框，
  % 使 6 字书名+卷名放进 9 字高的 chapter 框后不再压线（2026-10-05）
  shukou title chars = 4,
}

\title{史记}[梁孝王世家]
\author{汉·司马迁}
\maker{闲叟漫记·重製}


% ===== 索隐述赞样式（2026-11-08 四次定稿 · 愚大人亲手改定参数）=====
% ⚠ 竖排前提：自定义字体家族必须带 RawFeature={vertical:+vert}；
%   simhei.ttf 加 +vhal 会整体加载失败静默回退，只写 +vert；simfang 同理。
\newCJKfontfamily{\shuzanfont}[Path=C:/Windows/Fonts/,RawFeature={vertical:+vert}]{simhei.ttf}
\newCJKfontfamily{\shuzanbody}[Path=C:/Windows/Fonts/,RawFeature={vertical:+vert}]{simfang.ttf}
% 题头：152x30pt 圆角 6pt 棕底白字；baseline (0,2) 使整框落入界行格子不压版框线；
% 白字节点 x 不补偿、y +8.2pt（像素级墨迹实测 dx/dy≈0，L/R 留白约 20pt）。
\newcommand{\shuzanhead}[1]{%
  \begin{tikzpicture}[baseline={(0,2)}]
    \fill[sealbrown, rounded corners=6pt]
      (-76pt,-15pt) rectangle (76pt,15pt);
    \node[text=white, inner sep=0pt, anchor=center] at (0pt,8.2pt)
      {\shuzanfont\fontsize{23}{23}\selectfont #1};
  \end{tikzpicture}%
}
% ★ 修正书口卷名位置：chapter 框起始位置改回 cls 默认值，避免全局影响。
%   只在每个卷首页通过 \juan 命令里的空命令微调该页位置（安全，不会跨页污染）。
\ExplSyntaxOn
\AtBeginDocument{
  \dim_set:Nn \l__zhvt_shukou_chapter_title_start_dim
    { \l__zhvt_shukou_title_height_dim + \l__zhvt_font_size_dim * 2 }
  \dim_set:Nn \l__zhvt_shukou_page_number_start_dim
    { \l__zhvt_shukou_chapter_title_start_dim + \l__zhvt_shukou_chapter_title_height_dim }
}
\ExplSyntaxOff

\begin{document}

% ★★ 版心（书口）排布覆写：把书名放到鱼尾之下（2026-10-04）
%
%   cls 默认版心自上而下：书名（title 静态框，在鱼尾之上）→ 鱼尾 → 篇名（chapter 动态框）→ 页码
%   本模板改为：          鱼尾（位置不动）→ 书名 → 篇名 → 页码分割线 → 页码
%
%   手法：把 cls 放在 title 静态框里的书名，改放进 chapter 动态框（鱼尾之下）；
%         title 静态框留空。书名取 \l__zhvt_title_str，即 \title{...} 的第一参数，
%         所以换书时只改 \title，不必在此重复写书名。
%
%   ⚠ 三条必须遵守的写法（2026-10-04 实测）：
%     ① 汉字必须用 \hbox{...} 包住。flowfram 会把框内容装进一个宽 14.4mm（= 一个 line_sep
%        = 书口栏宽）的纵向 minipage 并按该宽度折行，而整页靠 PDF /Rotate 90 显示、
%        框内文字本身不变向 —— 不包 \hbox 时「围炉夜话」4 字会被逐字折成 4 行，
%        在版面上就成了横排 4 字并横向溢出版框。
%     ② 竖排下**一个 \hbox 就是一整列**。书名与篇名要同处一列，必须放进同一个 \hbox；
%        分成两个 \hbox 时，后一个会被当成新的一列、甩到隔壁栏去
%        （实测「序」跑到书口栏外侧 14.4mm 的正文列上）。中间的 \quad 在框内是横向空白，
%        显示出来正好是书名与篇名之间的纵向间隔。
%     ③ ★ 内容盒要写成 **与框等宽的定宽盒** `\hbox to \l__zhvt_shukou_chapter_title_height_dim{…\hfil}`。
%        实测：chapter 框是把内容盒沿纵向**居中**摆放的，内容宽度一变，书名就跟着漂：
%        自然宽度的 \hbox（书名 38.4mm）落在 x=157.2mm，加篇名后（57.6mm）落在 x=152.4mm，
%        目录页（无篇名）与正文页的书名就差半格。做成等宽盒后，无论有没有篇名，
%        内容固定从框顶起排，实测两页均落在 x=142.8mm，不再漂。
\ExplSyntaxOn
% 把书口 chapter 框（书名+篇名）起始位置下移 1 字高，避免 4 字卷名顶到鱼尾
% ⚠ cls 的 \__zhvt_calculate_dimensions: 在 \AtBeginDocument 里执行，会重算该 dim；
%   必须把覆写挂到 begin document 之后，否则会被二次覆盖（2026-10-05 实测）。
% (旧方案已移到导言区，直接修改 chapter 框起始位置)
% 卷标题命令 #1=书口篇名(鱼尾) #2=正文卷名（含卷序，卷名内空格用 ~）
%
% ★★ 目录页两项修订（2026-10-05，愚大人指出）★★
%   ① 页码起点统一对齐：
%      原状 —— 目录项写成 \par\hspace{4em}卷名\hspace{4em}页码，页码起点 = 卷名实宽 + 4em。
%      卷名实宽随字数变：「卷一~梁惠王上」6 字 +~、「卷七~离娄上」5 字 +~，
%      差 1 个字位 → 卷七~卷十 4 条的页码整体比其余 10 条高 27.21pt（一个字），看着上下不齐。
%      改法 —— 卷名补全角空格（宽 = 1 字 = 27.215pt）补足到 **7 字符位**再写进目录，
%      14 条页码起点遂全部落在同一条竖线上（起点对齐；末字随页码字数自然参差）。
%      ⚠ 只对目录里的副本 \l_tmpc_tl 补，正文卷名仍用 \l_tmpa_tl 原样 —— 正文像素不受影响。
%   ② 卷名整体上移半字位：前导 \hspace 由 4em 改 3.5em。
%      实测原状目录块在版框内偏下（天头留白 139.6pt / 地脚留白 110.7pt），
%      上移 0.5em = 13.6pt 后上下基本均分，视觉居中。
\NewDocumentCommand { \juan } { mm } {
	\clearpage
	\tl_set:Nn \l__zhvt_chapter_title_tl { #1 }
	\tl_set:Nn \l_tmpa_tl { #2 }
	\tl_set_eq:NN \l_tmpc_tl \l_tmpa_tl
	\int_step_inline:nn { 7 - \str_count:n { #2 } }
	  { \tl_put_right:Nn \l_tmpc_tl { 　 } }
	\tl_set:Nx \l_tmpb_tl { \zhnum { page } }
	\tl_put_left:Nn \l_tmpb_tl { \hspace { 4em } }
	\tl_put_left:Nn \l_tmpb_tl { \l_tmpc_tl }
	\tl_put_left:Nn \l_tmpb_tl { \par \hspace { 3.5em } }
	\addtocontents { toc } { \l_tmpb_tl }
	\skip_horizontal:n { \l__zhvt_font_size_dim * 4 }
	\l_tmpa_tl
	\par
	\thispagestyle { empty }
}
\cs_set_protected:Npn \__zhvt_set_shukou_content: {
	% ★ 鱼尾上方留白（title 静态框）嵌入闲叟钤印：
	%   整页 /Rotate 90，故此处把图用 \rotatebox{90} 转正为竖章；
	%   尺寸卡在版心条幅内（书内约 7mm 宽 × 33.8mm 高，留余量不压鱼尾）。
	\setstaticcontents * { title } {
		\vbox_to_ht:nn { \l__zhvt_baseline_skip_dim } {
			\vskip 14.2pt
			\hbox_to_wd:nn { \l__zhvt_shukou_title_height_dim } {
				\hfil
				\rotatebox [ origin = c ] { 90 } { \includegraphics [ width = 7mm , height = 33.8mm ] { yshs.png } }
				\hfil
			}
		}
	}
% ★ 版心（书口）统一改法（2026-11-16 愚大人定）：18pt＋框内居中＋书口文字同版框色
%   ① 原 27.2pt 时 7 字长篇名（如「廉颇蔺相如列传」）末字溢出 chapter 框、压住页码；
%      改 \fontsize{18pt}{18pt}，前后 \hfil 令字串在框内居中（不随篇名字数漂移）。
%   ② \vbox{\vskip 8.2pt} 定量下移，抵消缩字号后栏宽方向 8.2pt 上偏（像素实测 y 中心 543.3 = 原版）。
%      ⚠ \raisebox / \vspace* 在 flowfram 框内无效；换字号须同步换补偿：16pt→10.1、19pt→7.4。
%   ③ chapter 与 page~number 均置 \color{sealbrown}（HTML 874434），与版框／界行同色；正文仍黑。
%   ⚠ 页码字位、钤印、鱼尾、界行、正文一律未动；五部合订本同此口径。
	\setdynamiccontents * { chapter }
	{ \color{sealbrown} \vbox { \vskip 8.2pt \hbox to \l__zhvt_shukou_chapter_title_height_dim
		{ \hfil \fontsize{18pt}{18pt}\selectfont \l__zhvt_title_str \quad \l__zhvt_chapter_title_tl \hfil } } }
	\setdynamiccontents * { page~number } { \color{sealbrown} \quad \zhnum { page } }
}
\__zhvt_set_shukou_content:
\ExplSyntaxOff

\maketitle

% 单篇，无多卷目录，不渲染目录页

\juan{梁孝王世家}{卷五十八~梁孝王世家第二十八}
\hspace*{2em}梁孝王武者。孝文皇帝子也。而与孝景帝同母。母。窦太后也。\\
\hspace*{2em}孝文帝凡四男。长子曰太子。是为孝景帝。次子武。次子参。次子胜。孝文帝即位二年。以武为代王。以参为太原王。以胜为梁王。二岁。徙代王为淮阳王。以代尽与太原王。号曰代王。参立十七年。孝文後二年卒。谥为孝王。子登嗣立。是为代共王。立二十九年。元光二年卒。子义立。是为代王。十九年。汉广关。以常山为限。而徙代王王清河。清河王徙以元鼎三年也。\\
\hspace*{2em}初。武为淮阳王十年。而梁王胜卒。谥为梁怀王。怀王最少子。爱幸异於他子。其明年。徙淮阳王武为梁王。梁王之初王梁。孝文帝之十二年也。梁王自初王通历已十一年矣。\\
\hspace*{2em}梁王十四年。入朝。十七年。十八年。比年入朝。留。其明年。乃之国。二十一年。入朝。二十二年。孝文帝崩。二十四年。入朝。二十五年。复入朝。是时上未置太子也。上与梁王燕饮。尝从容言曰。﹁千秋万岁後传於王。﹂王辞谢。虽知非至言。然心内喜。太后亦然。\\
\hspace*{2em}其春。吴楚齐赵七国反。吴楚先击梁棘壁。杀数万人。梁孝王城守睢阳。而使韩安国。张羽等为大将军。以距吴楚。吴楚以梁为限。不敢过而西。与太尉亚夫等相距三月。吴楚破。而梁所破杀虏略与汉中分。明年。汉立太子。其後梁最亲。有功。又为大国。居天下膏腴地。地北界泰山。西至高阳。四十馀城。皆多大县。\\
\hspace*{2em}孝王。窦太后少子也。爱之。赏赐不可胜道。於是孝王筑东苑。方三百馀里。广睢阳城七十里。大治宫室。为衤复道。自宫连属於平台三十馀里。得赐天子旌旗。出从千乘万骑。东西驰猎。拟於天子。出言。入言警。招延四方豪桀。自山以东游说之士。莫不毕至。齐人羊胜。公孙诡。邹阳之属。公孙诡多奇邪计。初见王。赐千金。官至中尉。梁号之曰公孙将军。梁多作兵器弩弓矛数十万。而府库金钱且百巨万。珠玉宝器多於京师。\\
\hspace*{2em}二十九年十月。梁孝王入朝。景帝使使持节乘舆驷马。迎梁王於关下。既朝。上疏因留。以太后亲故。王入则侍景帝同辇。出则同车游猎。射禽兽上林中。梁之侍中。郎。谒者著籍引出入天子殿门。与汉宦官无异。\\
\hspace*{2em}十一月。上废栗太子。窦太后心欲以孝王为後嗣。大臣及袁盎等有所关说於景帝。窦太后义格。亦遂不复言以梁王为嗣事由此。以事秘。世莫知。乃辞归国。\\
\hspace*{2em}其夏四月。上立胶东王为太子。梁王怨袁盎及议臣。乃与羊胜。公孙诡之属阴使人刺杀袁盎及他议臣十馀人。逐其贼。未得也。於是天子意梁王。逐贼。果梁使之。乃遣使冠盖相望於道。覆按梁。捕公孙诡。羊胜。公孙诡。羊胜匿王後宫。使者责二千石急。梁相轩丘豹及内史韩安国进谏王。王乃令胜。诡皆自杀。出之。上由此怨望於梁王。梁王恐。乃使韩安国因长公主谢罪太后。然后得释。\\
\hspace*{2em}上怒稍解。因上书请朝。既至关。茅兰说王。使乘布车。从两骑入。匿於长公主园。汉使使迎王。王已入关。车骑尽居外。不知王处。太后泣曰。﹁帝杀吾子。﹂景帝忧恐。於是梁王伏斧质於阙下。谢罪。然後太后。景帝大喜。相泣。复如故。悉召王从官入关。然景帝益疏王。不同车辇矣。\\
\hspace*{2em}三十五年冬。复朝。上疏欲留。上弗许。归国。意忽忽不乐。北猎良山。有献牛。足出背上。孝王恶之。六月中。病热。六日卒。谥曰孝王。\\
\hspace*{2em}孝王慈孝。每闻太后病。口不能食。居不安寝。常欲留长安侍太后。太后亦爱之。及闻梁王薨。窦太后哭极哀。不食。曰。﹁帝果杀吾子。﹂景帝哀惧。不知所为。与长公主计之。乃分梁为五国。尽立孝王男五人为王。女五人皆食汤沐邑。於是奏之太后。太后乃说。为帝加壹餐。\\
\hspace*{2em}梁孝王长子买为梁王。是为共王。子明为济川王。子彭离为济东王。子定为山阳王。子不识为济阴王。\\
\hspace*{2em}孝王未死时。财以巨万计。不可胜数。及死。藏府馀黄金尚四十馀万斤。他财物称是。\\
\hspace*{2em}梁共王三年。景帝崩。共王立七年卒。子襄立。是为平王。\\
\hspace*{2em}梁平王襄十四年。母曰陈太后。共王母曰李太后。李太后。亲平王之大母也。而平王之后姓任。曰任王后。任王后甚有宠於平王襄。初。孝王在时。有罍樽。直千金。孝王诫後世。善保罍樽。无得以与人。任王后闻而欲得罍樽。平王大母李太后曰。﹁先王有命。无得以罍樽与人。他物虽百巨万。犹自恣也。﹂任王后绝欲得之。平王襄直使人开府取罍樽。赐任王后。李太后大怒。汉使者来。欲自言。平王襄及任王后遮止。闭门。李太后与争门。措指。遂不得见汉使者。李太后亦私与食官长及郎中尹霸等士通乱。而王与任王后以此使人风止李太后。李太后内有淫行。亦已。後病薨。病时。任后未尝请病。薨。又不持丧。\\
\hspace*{2em}元朔中。睢阳人类犴反者。人有辱其父。而与淮阳太守客出同车。太守客出下车。类犴反杀其仇於车上而去。淮阳太守怒。以让梁二千石。二千石以下求反甚急。执反亲戚。反知国阴事。乃上变事。具告知王与大母争樽状。时丞相以下见知之。欲以伤梁长吏。其书闻天子。天子下吏验问。有之。公卿请废襄为庶人。天子曰。﹁李太后有淫行。而梁王襄无良师傅。故陷不义。﹂乃削梁八城。枭任王后首于市。梁馀尚有十城。襄立三十九年卒。谥为平王。子无伤立为梁王也。\\
\hspace*{2em}济川王明者。梁孝王子。以桓邑侯孝景中六年为济川王。七岁。坐射杀其中尉。汉有司请诛。天子弗忍诛。废明为庶人。迁房陵。地入于汉为郡。\\
\hspace*{2em}济东王彭离者。梁孝王子。以孝景中六年为济东王。二十九年。彭离骄悍。无人君礼。昏暮私与其奴。亡命少年数十人行剽杀人。取财物以为好。所杀发觉者百馀人。国皆知之。莫敢夜行。所杀者子上书言。汉有司请诛。上不忍。废以为庶人。迁上庸。地入于汉。为大河郡。\\
\hspace*{2em}山阳哀王定者。梁孝王子。以孝景中六年为山阳王。九年卒。无子。国除。地入于汉。为山阳郡。\\
\hspace*{2em}济阴哀王不识者。梁孝王子。以孝景中六年为济阴王。一岁卒。无子。国除。地入于汉。为济阴郡。\\
\hspace*{2em}太史公曰。梁孝王虽以亲爱之故。王膏腴之地。然会汉家隆盛。百姓殷富。故能植其财货。广宫室。车服拟於天子。然亦僭矣。\\
\hspace*{2em}褚先生曰。臣为郎时。闻之於宫殿中老郎吏好事者称道之也。窃以为令梁孝王怨望。欲为不善者。事从中生。今太后。女主也。以爱少子故。欲令梁王为太子。大臣不时正言其不可状。阿意治小。私说意以受赏赐。非忠臣也。齐如魏其侯窦婴之正言也。何以有後祸。景帝与王燕见。侍太后饮。景帝曰。﹁千秋万岁之後传王。﹂太后喜说。窦婴在前。据地言曰。﹁汉法之约。传子適孙。今帝何以得传弟。擅乱高帝约乎。﹂於是景帝默然无声。太后意不说。\\
\hspace*{2em}故成王与小弱弟立树下。取一桐叶以与之。曰。﹁吾用封汝。﹂周公闻之。进见曰。﹁天王封弟。甚善。﹂成王曰。﹁吾直与戏耳。﹂周公曰。﹁人主无过举。不当有戏言。言之必行之。﹂於是乃封小弟以应县。是後成王没齿不敢有戏言。言必行之。孝经曰。﹁非法不言。非道不行。﹂此圣人之法言也。今主上不宜出好言於梁王。梁王上有太后之重。骄蹇日久。数闻景帝好言。千秋万世之後传王。而实不行。\\
\hspace*{2em}又诸侯王朝见天子。汉法凡当四见耳。始到。入小见。到正月朔旦。奉皮荐璧玉贺正月。法见。後三日。为王置酒。赐金钱财物。後二日。复入小见。辞去。凡留长安不过二十日。小见者。燕见於禁门内。饮於省中。非士人所得入也。今梁王西朝。因留。且半岁。入与人主同辇。出与同车。示风以大言而实不与。令出怨言。谋畔逆。乃随而忧之。不亦远乎。非大贤人。不知退让。今汉之仪法。朝见贺正月者。常一王与四侯俱朝见。十馀岁一至。今梁王常比年入朝见。久留。鄙语曰﹁骄子不孝﹂。非恶言也。故诸侯王当为置良师傅。相忠言之士。如汲黯。韩长孺等。敢直言极谏。安得有患害。\\
\hspace*{2em}盖闻梁王西入朝。谒窦太后。燕见。与景帝俱侍坐於太后前。语言私说。太后谓帝曰。﹁吾闻殷道亲亲。周道尊尊。其义一也。安车大驾。用梁孝王为寄。﹂景帝跪席举身曰。﹁诺。﹂罢酒出。帝召袁盎诸大臣通经术者曰。﹁太后言如是。何谓也。﹂皆对曰。﹁太后意欲立梁王为帝太子。﹂帝问其状。袁盎等曰。﹁殷道亲亲者。立弟。周道尊尊者。立子。殷道质。质者法天。亲其所亲。故立弟。周道文。文者法地。尊者敬也。敬其本始。故立长子。周道。太子死。立適孙。殷道。太子死。立其弟。﹂帝曰。﹁於公何如。﹂皆对曰。﹁方今汉家法周。周道不得立弟。当立子。故春秋所以非宋宣公。宋宣公死。不立子而与弟。弟受国死。复反之与兄之子。弟之子争之。以为我当代父後。即刺杀兄子。以故国乱。祸不绝。故春秋曰﹁君子大居正。宋之祸宣公为之﹂。臣请见太后白之。﹂袁盎等入见太后。﹁太后言欲立梁王。梁王即终。欲谁立。﹂太后曰。﹁吾复立帝子。﹂袁盎等以宋宣公不立正。生祸。祸乱後五世不绝。小不忍害大义状报太后。太后乃解说。即使梁王归就国。而梁王闻其义出於袁盎诸大臣所。怨望。使人来杀袁盎。袁盎顾之曰。﹁我所谓袁将军者也。公得毋误乎。﹂刺者曰。﹁是矣。﹂刺之。置其剑。剑著身。视其剑。新治。问长安中削厉工。工曰。﹁梁郎某子来治此剑。﹂以此知而发觉之。发使者捕逐之。独梁王所欲杀大臣十馀人。文吏穷本之。谋反端颇见。太后不食。日夜泣不止。景帝甚忧之。问公卿大臣。大臣以为遣经术吏往治之。乃可解。於是遣田叔。吕季主往治之。此二人皆通经术。知大礼。来还。至霸昌厩。取火悉烧梁之反辞。但空手来对景帝。景帝曰。﹁何如。﹂对曰。﹁言梁王不知也。造为之者。独其幸臣羊胜。公孙诡之属为之耳。谨以伏诛死。梁王无恙也。﹂景帝喜说。曰。﹁急趋谒太后。﹂太后闻之。立起坐餐。气平复。故曰。不通经术知古今之大礼。可以为三公及左右近臣。少见之人。如从管中闚天也。\\

% ===== 索隐述赞（司马贞《史记索隐》）· 统一样式 =====
% 题头「索隐述赞曰」五字并入反白框；正文每四句一列，定宽盒 \hfil 居列中；
% \fontsize{24}{41} 列距对齐全书界行（121+41k）；\clearpage 使述赞独占一页。
\clearpage
\shuzanhead{索隐述赞曰}\linebreak
{\shuzanbody\fontsize{24}{41}\selectfont
\hbox to \hsize{\hfil 文帝少子。徙封於梁。太后锺爱。广筑睢阳。\hfil}\linebreak
\hbox to \hsize{\hfil 旌旂警跸。势拟天王。功扞吴楚。计丑孙羊。\hfil}\linebreak
\hbox to \hsize{\hfil 窦婴正议。袁盎劫伤。汉穷梁狱。冠盖相望。\hfil}\linebreak
\hbox to \hsize{\hfil 祸成骄子。致此猖狂。虽分五国。卒亦不昌。\hfil}%
\par}
\end{document}
